\section*{Tóm tắt}
\thispagestyle{empty}

Trong những năm gần đây, thể loại game thẻ bài chiến thuật (Strategy Card Game) tiếp tục thu hút đông đảo người chơi nhờ yêu cầu cao về tư duy, lập kế hoạch và khả năng ra quyết định. Tuy nhiên, việc phát triển các trò chơi thẻ bài đối kháng trực tuyến đa người chơi đặt ra nhiều thách thức về thiết kế kiến trúc phần mềm, đặc biệt là trong việc quản lý vòng đời đối tượng, đồng bộ trạng thái mạng, kiểm soát quyền hạn và khả năng mở rộng hệ thống.

Đề tài này tập trung xây dựng một hệ thống game thẻ bài đối kháng nhiều người chơi trên nền tảng Unity, với kiến trúc Unity Client được thiết kế theo hướng tiến hóa từ mô hình MVC truyền thống sang kiến trúc hướng subsystem. Trong kiến trúc này, Controller đóng vai trò như các \textit{application service}, chịu trách nhiệm xử lý nghiệp vụ và kiểm soát toàn bộ việc thay đổi trạng thái, trong khi View chỉ đảm nhiệm vai trò hiển thị và phát sinh tương tác người dùng. Vòng đời của các subsystem được quản lý thống nhất thông qua Dependency Injection (Zenject), giúp hệ thống linh hoạt trước vòng đời động của Unity và giảm thiểu coupling giữa các thành phần.

Hệ thống mạng trong Unity Client được hiện thực dựa trên Photon Fusion 2 với cơ chế đồng bộ trạng thái (state-based synchronization). Các Model mạng được đồng bộ thông qua các thuộc tính \texttt{[Network]}, trong khi Controller đảm bảo mô hình multiplayer authority, ngăn chặn các cập nhật trạng thái trái phép từ client không có thẩm quyền. Bên cạnh đó, kiến trúc hỗ trợ cơ chế replay thông qua việc ghi nhận các hành động gameplay quan trọng dưới dạng Command, được tạo và ghi log nội bộ trong Controller, phục vụ tái hiện trận đấu và phân tích sau trận.

Song song với Unity Client, hệ thống Backend được xây dựng bằng ASP.NET Core theo kiến trúc phân tầng đóng (layered architecture). Backend chịu trách nhiệm quản lý dữ liệu người dùng, thông tin thẻ bài, lịch sử đấu và các chức năng nghiệp vụ ngoài thời gian thực. Các API Backend cung cấp giao diện thống nhất để Unity Client truy cập dữ liệu, trong khi cơ chế xác thực và phân quyền đảm bảo an toàn hệ thống.

Đồ án bao gồm các nội dung chính: phân tích yêu cầu và các mô hình kiến trúc game hiện đại, thiết kế kiến trúc Unity Client hướng subsystem kết hợp Dependency Injection, hiện thực hệ thống mạng bằng Photon Fusion 2, xây dựng cơ chế replay và thiết kế Backend ASP.NET Core hỗ trợ quản lý dữ liệu. Kết quả đạt được là một hệ thống game thẻ bài đối kháng hoạt động ổn định, có khả năng mở rộng, dễ bảo trì và phù hợp cho các kịch bản multiplayer hiện đại, đồng thời đóng vai trò như một tài liệu tham khảo cho việc thiết kế kiến trúc game trong môi trường Unity.


\clearpage
\pagenumbering{arabic}